\section{Supporting pore-flow, rock and older-kernel timing evidence}
\subsection{Analytic pore flux and capillary entry}
Two parallel, isolated slit pores have widths $H_0=8s$ and $H_1=16s$,
with total specimen height $36s$, for $s=1,2,3$. Each pore contains a pure
component; separate solid walls prevent a fluid-fluid interface. Periodic
flow is driven by $F_x=10^{-6}s^{-3}$ with both permutations of
$(\tau_0,\tau_1)=(0.6,1.4)$. The independently sampled Poiseuille reference is
\begin{equation}
 q_k=\frac{F_x}{\nu_k(36s)}\left(\frac{H_k^3}{12}+\frac{H_k}{24}\right).
\end{equation}
The second term accounts for midpoint sampling of the analytic velocity.
Single-phase runs on the same geometry supply $K$ for normalization. Across
both D2Q9 and extruded D3Q19, the maximum phase-flux discrepancy is
$\BundleFluxError$ (Fig.~\ref{fig:pores}), consistent with the remaining
transient at the selected duration and roundoff. This verifies phase attribution and
viscosity normalization in a controlled pore system; it is not evidence for
general relative-permeability curves in connected rock pores.

Pressure boundaries then drive phase 0 into an initially phase-1-filled
60-degree slit. Its sharp static reference is
$p_c=2\sigma\cos\theta/H$. At pressure drops $0.5p_c$ and $1.5p_c$, the
late interface respectively retreats and advances. Diffusive refinement uses
$H=16s$, $\sigma=0.005/s$ and duration $8000s^2$, for $s=1,2,4$.
Four D3Q19 square-tube controls at $H=16$ use the force-balance reference
$p_c=4\sigma\cos\theta/H$. Both gradient variants give the expected late
direction in these controls. In the square tube, however, inferring an
entry pressure from a developed-flow hydraulic resistance gives deviations
of approximately 17--25\% for the clipped stencil and 9--16\% for
renormalization. Contact-line dissipation, voxel corners and diffuse wetting
are not removed by this diagnostic. The result checks the direction of
displacement and quantifies limitations; it does not calibrate a static
capillary-pressure curve or a dynamic contact-angle law.

\begin{figure}[htbp]\centering
\includegraphics[width=.95\linewidth]{pore-controls.pdf}
\caption{Independent slit-bundle phase-flux checks and pressure-driven
capillary-entry controls. Hydraulic pressure differences are diagnostic
deviations from a sharp static reference, not a calibrated contact-line model.\AIPlotDisclosure}
\label{fig:pores}\end{figure}

\subsection{Public Bentheimer specimen}
The segmented 4.54~$\mu$m Bentheimer image of Shah, Crawshaw and Boek
\cite{shah2014} supplies a public, CC BY 4.0 input. A centered
$48^3$ crop starts at $(263,263,263)$ in the original $575^3$ array.
Its porosity is $\RockPorosity$, and all pores belong to one connected
component after periodic face connections. No pore removal, filtering or
segmentation change is performed. Phase 0 initially occupies pores whose
periodic distance to the grain-center mask is at least three original
voxels; phase 1 occupies the remaining pore space. This deterministic
morphological rule specifies an initial state and does not prescribe an
experimental drainage history.

Single-phase MRT flow at $\tau=1$ yields
$K=\RockK\ \mathrm{m}^2$. Halving and doubling the force give a relative
spread of $\RockForceSpread$. Independent LBPM v2026.04 calculations first
pass three analytic slit controls, then use the same exact periodically
padded mask, force, viscosity and averaging volume. Across five matched
rock cases, the largest permeability difference is
\RockReferenceDifference\ percent. Both solvers retain small finite-density
convention differences; converted library permeability columns are not used.
Pure-phase color-flow reductions agree with density-normalized SCSP velocity
within $10^{-11}$ in the tested cases. These are numerical crosschecks on
the specified periodically repeated specimen, not a comparison with
laboratory sandstone permeability.

Each original voxel is also divided into eight voxels, preserving the
binary box geometry and initial phase volumes. Force, surface tension and
duration are scaled diffusively. The resulting single-phase permeability
decreases by \RefinedKChange\ percent. The independent refined LBPM reference
reproduces this shift. This is a discretization sensitivity
on a fixed segmented specimen; two resolutions do not establish an
asymptotic convergence rate, and subdivision adds no imaging information.

Closed two-phase runs at prescribed angles 60 and 120 degrees retain
component masses to a maximum relative drift of $\RockMassDrift$.
The zero-force control also evolves, so its motion cannot be attributed to
the imposed body force. Physical interface relaxation and possible numerical
currents have not been separated in this geometry. In the forced 60-degree case, the color
field has RMS change \RockLateColor\ over the last 4000 steps of a
16000-step run. The subdivided case is compared at the matched base time
8000, not equated with the longer coarse endpoint. Planar controls at the
two rock-study surface tensions measure a bulk interface width of 5.289
lattice nodes, while the original crop's median periodic
distance to grain centers is only two voxels. These lengths are not a
calibrated throat-radius measurement, but they expose inadequate scale
separation in many pores. Transient phase fluxes in Fig.~\ref{fig:rock} are
therefore retained as measured observables and are not promoted to steady
relative permeability or experimental wetting predictions.

\begin{figure}[htbp]\centering
\includegraphics[width=.95\linewidth]{digital-rock.pdf}
\caption{Actual phase fields and advective flux histories on the attributed
Bentheimer crop. Grain voxels are masked. The zero-force control and
late field changes distinguish evolving unforced motion from a steady
transport measurement. The refined curve is scaled to the original
diffusive units and is compared at matching times.\AIPlotDisclosure}
\label{fig:rock}\end{figure}

Two additional full-field comparisons on this rock are bitwise identical:
one versus two MPI ranks, and uninterrupted versus a one-to-two-rank restart.
These extend the homogeneous/droplet replay checks to an irregular pore mask.

\subsection{Workstation performance}
This timing campaign uses the frozen perturbation executable recorded before
the experimental capillary extension. It does not measure the new CSF branch
or establish the runtime of the revised build. The complete executable hash
remains in the timing data.
A separate $96^3$ crop provides a fixed two-phase timing workload with
\PerfFluidNodes\ fluid nodes. The actual D3Q19/MRT numerical path is timed
for \PerfSteps\ updates at one, two, four and eight MPI ranks on an
Intel Core Ultra 9 285K workstation. One full warm-up per rank count is
discarded; three sequential measured repetitions interleave rank counts.
Each local rank is pinned to one of the recorded logical CPUs. Frequency
is governed by the operating system, not fixed by the experiment.
All reconstructed fields agree bitwise across ranks and repetitions.

The numerical timer includes streaming, MPI exchange, boundaries,
macroscopic reconstruction, collision and finite-state checks. Initial
collision is initialization, not an extra counted update. Separate timers
cover field/diagnostic output and population checkpoints; construction and
launcher overhead belong to the process duration. Each timer is a maximum
across ranks, so category maxima need not sum to the maximum run wall time.
Throughput counts fluid-node updates, although dense solid arrays also incur
streaming and memory costs. Fig.~\ref{fig:performance} reports median and
observed range, with exact counts and per-rank memory in the supplement.
These are local CPU/MPI measurements; they establish neither cluster
scaling nor GPU performance, and no runtime comparison with LBPM is inferred.

\begin{figure}[htbp]\centering
\includegraphics[width=.95\linewidth]{workstation-performance.pdf}
\caption{Repeated, fixed-affinity workstation timing. Error bars show the
observed minimum and maximum of three runs. Numerical work and I/O are
reported separately; the speedup denominator is the same measured workload
at one rank.\AIPlotDisclosure}
\label{fig:performance}\end{figure}

